\section{Discussion}
\label{sec:discussion}

The protocol is intended for evaluations where the performance object is broader than final
accuracy: multi-provider API studies, candidate-pool and selector studies, router or cascade
studies, closed commercial APIs, incomplete telemetry, failed or unscorable cells, and cost- or
latency-sensitive inference. In those settings, the central practical lesson is to report the layer
at which each claim is measured. A nominal call budget, the realized resource vector, the cost of
constructing a candidate pool, and the marginal cost of a selector are different performance
objects.

\paragraph{Compare selectors on fixed candidate pools}
When a study evaluates a nontrivial commit rule, report accuracy against plurality on the same
extracted answers and disclose the pool-construction cost. This separates selection from discovery
and prevents generator-relative gains from being mistaken for selector progress
(Sections~\ref{sec:identical-pool} and~\ref{sec:heldout-calibrated-selector}).

\paragraph{Treat failure-mined rules as transfer claims}
Rules developed from one provider or workload should carry a development-environment label and be
reported with provider-stratified rescue/regression counts. The relevant question is not only
whether the rule helps somewhere, but whether its failure signature transfers under the evaluated
provider shifts (Section~\ref{sec:provider-analysis}).

\paragraph{Report nominal budget with componentwise realized resources}
Equal nominal budgets do not by themselves equalize successful completions, retries, generated
tokens, latency, or dollars. Resource fields should therefore be reported componentwise, with
reconstruction status and lower-bound fields marked explicitly (Section~\ref{sec:compute}).

\paragraph{Preserve blocked outcomes and discovery--selection structure}
Protocol-blocked cells should remain visible rather than being scored as zero or silently removed.
For completed cells, joint reporting of discovery rate and conditional selection accuracy identifies
whether errors are absent-from-pool failures or present-but-misselected failures
(Sections~\ref{sec:results} and~\ref{sec:setup}).

These practices are actionable for future closed-API evaluations because they do not require access
to provider internals. They require disciplined accounting: the evaluation must say which resources
were observed, which were reconstructed, which remain lower bounds, which cells failed the bounded
protocol, and which selector comparisons share the same candidate pool.
